% ============================================================================

% Paper: Frequent Losers as Cover Bidders in Public Procurement
% ============================================================================

\section{Empirical Strategy}
\label{sec:empirical}

\paragraph{Estimand.}
We estimate $\beta$, the conditional mean difference in log
negotiated price between tenders that include at least one FL
firm and otherwise comparable tenders, within item type, year,
and purchasing unit. We are explicit that $\beta$ is a
\emph{conditional association}, not a causal parameter unless
the identification assumptions stated below are satisfied. We
identify $\beta$ under two complementary designs of unequal
strength: a baseline conditional-on-observables specification
(\S\ref{sec:reduced_form}), which provides the descriptive
range; and an institutional design exploiting the variation
the convite minimum-bidder rule creates
(\S\ref{sec:additional_id}), which is our primary
identification. The institutional design generates three
sources of variation that interact:

\begin{enumerate}
\item \textbf{Modal comparison} (convite vs.\ pregão).
Convite (Lei \valLeiOriginal Art.\ 22) requires three valid proposals;
pregão (Lei \valLeiPregao) has no analogous rule. Comparing
FL-present tenders across modalities within the same buyer,
item-group, and year identifies the regulatory cover-bidding
channel under the assumption that modal assignment is
orthogonal to FL classification conditional on item-year-buyer
fixed effects.
\item \textbf{Within-convite constraint binding}. Within the
convite subsample, the rule binds when $n_{\text{genuine}} < 3$
(cover bidding statutorily required) and slack when
$n_{\text{genuine}} \geq 3$ (cover bidding voluntary). Under
the assumption that the binding-versus-slack split is
orthogonal to the unobserved cartel-selection process
conditional on observed market characteristics---we discuss
sensitivity to this assumption below---the
voluntary subset identifies the strategic component of FL
deployment, separated from mechanical procedural compliance.
\item \textbf{Statutory value threshold}. Lei \valLeiOriginal Art.\ 23
caps convite procurement at \valCapOldR (\valCapNewR after
Decreto \valDecreto); above the cap, only pregão is permitted.
The cap is administered on documented contract value, not
bidder identity. We use the threshold not as a fuzzy RDD---
local first-stage power is too weak (Stock--Yogo $F$ below
threshold)---but as an exogeneity check on modal assignment.
McCrary-style density tests around the \valCapOldR cap find no
evidence of bunching (left/right density ratio $\valMcCraryRatio$,
log-density discontinuity $\valMcCraryDisc$;
Table~\ref{tab:mccrary}); the convite share drops 9.8
percentage points across the cap (\valPctSixnine below, \valPctFifnine above),
consistent with rule-driven assignment rather than strategic
gaming. The threshold check rules out the strongest version of
the alternative that buyers self-select into convite to deploy
cover bidders.
\end{enumerate}

\paragraph{Identifying assumptions and what we test.}
The institutional design identifies $\beta$ under three
assumptions, each with a sensitivity analysis. \emph{(A.1)
Modal-assignment exogeneity}: convite vs.\ pregão modal choice
is uncorrelated with the unobserved cartel-selection process
conditional on item-year-buyer fixed effects. We test the
strongest violation---selection on contract value to deploy
cover bidders---via the McCrary density check at the \valCapOldR
cap, which passes. \emph{(A.2) Constraint-binding exogeneity}:
the within-convite split between binding and slack is
uncorrelated with unobserved cartel selection conditional on
the FE structure. We bound violations using the
\citet{cinelli2020making} sensitivity analysis, which yields
$RV_{q=1} = \valRVqOne$: an unobserved confounder would need to
explain at least \valRVqOne of residual variation in both FL
presence and log prices to nullify the coefficient. \emph{(A.3)
No anticipation}: market entry of FL firms is not driven by
contemporaneous price expectations within the same item-PBU
cell. The reverse-causality test in
Section~\ref{sec:mechanisms} (M3) shows the elasticity of FL
entry with respect to lagged prices is $\valMechMThreeCoef$ (SE $\valMechMThreeSE$);
price-chasing is too weak to account for the main results.

The three assumptions do not jointly establish strict causal
identification: pre-event coefficients in the within-PBU FL-entry
event study at $t = -2,-3$ remain positive and similar to the
post-event coefficient (Figure~\ref{fig:fl_entry_es},
\ref{app:robustness}). We interpret this pre-trend as evidence
of cartel selection \emph{into} markets that already exhibit
high prices---the very behavior the screen is designed to
flag---rather than as a violation of A.1--A.3. Under that
interpretation, the institutional design identifies the
\emph{premium that FL participation generates conditional on
cartel-selected markets}, not the average causal effect of FL
participation on prices in a randomly drawn market. We are
explicit about this in Section~\ref{sec:limitations}: the
identification we deliver is plausible and bounded by the
sensitivity analysis, not strict.

\paragraph{Outline.}
Section~\ref{sec:reduced_form} reports the OLS, cross-fit, and
matching baselines. Section~\ref{sec:detection_validation}
benchmarks the screen against CADE cartel co-bidding ground
truth. Sections~\ref{sec:supporting_setup}--\ref{sec:empirical_network}
present supporting diagnostics (structural estimation,
Bajari--Ye coordination tests, network-split heterogeneity) that
assess coherence with cover-bidding mechanisms.
Section~\ref{sec:additional_id}, embedded in the Results,
reports the institutional-identification estimates that the
present section sets up.

% ── 5.1 Reduced-Form Identification ──────────────────────────────

\subsection{Reduced-Form Conditional Associations}
\label{sec:reduced_form}

\noindent\textit{The institutional identification of
Section~\ref{sec:additional_id} is our primary design. The
specifications below are conditional-association baselines,
estimated on the same regression sample, that bracket the
modal-channel estimates and document the FL--price gap
descriptively.}\medskip

\paragraph{OLS baseline.}
The baseline specification regresses tender outcomes on FL presence:
\begin{equation}
  y_{igt} = \beta \cdot \text{losers}_{igt} + \mathbf{x}_{igt}'
  \boldsymbol{\delta} + \alpha_g + \lambda_t + \gamma_k + \varepsilon_{igt}
  \label{eq:ols_main}
\end{equation}
where $y_{igt}$ is the outcome for tender-item $i$ in item group $g$ at
time $t$ and purchasing unit $k$; $\text{losers}_{igt} \in \{0,1\}$
indicates FL presence; $\mathbf{x}_{igt}$ includes a convite indicator;
$\alpha_g$, $\lambda_t$, and $\gamma_k$ are item, year, and PBU fixed
effects; and $\varepsilon_{igt}$ is the error term clustered at the item
level.

The OLS coefficient $\hat{\beta}$ captures a conditional
association---the average difference in outcomes between FL-present and
FL-absent tenders within the same item type, year, and purchasing unit.
Under the cover-bidding hypothesis, $\hat{\beta} > 0$ for prices.
This gap is what one would expect if cover bidding inflates prices, but
it could also reflect unobserved market-level factors that attract both
FL firms and higher prices. Item, year, and PBU fixed effects absorb
time-invariant cross-market heterogeneity; the residual threat is time-varying within-market confounders---for
instance, demand shocks or the entry of large firms that
simultaneously raise prices and attract FL participation. We assess the sensitivity of
$\hat{\beta}$ to such confounders using the \citet{cinelli2020making}
framework, which yields a robustness value of $RV_{q=1} = \valRVqOne$: an
unobserved confounder would need to explain at least \valRVqOne of the
residual variation in both FL presence and log prices to reduce
$\hat{\beta}$ to zero.

\paragraph{Matching estimators.}
To address potential selection on observables, we estimate the
price effect using two matching approaches: coarsened exact
matching (CEM) on year, procedure type, item group, and PBU size
quartile ($\hat{\beta}_{\text{CEM}} = 0.077$, $N = \valCleanItems$), and
inverse probability weighting (IPW, $\hat{\beta}_{\text{IPW}} = 0.055$,
$N = \valFilteredItems$). These estimates bracket the OLS baseline (0.064)
and the cross-fit average (0.036), providing a non-parametric check
that the price association survives reweighting on observables.

\paragraph{Instrumental variable (measurement-error diagnostic).}
The binary FL indicator contains misclassification noise: firms near
the IQR threshold may not be cover bidders, and some cover bidders
escape the threshold. Classical binary measurement error attenuates
the OLS coefficient toward zero. We use a leave-one-out FL supply
instrument primarily to diagnose this attenuation:
\begin{equation}
  Z_{kgt} = \sum_{j \neq k} \mathbf{1}[\text{FL firm active at PBU } j
  \text{ in group } g, \text{ year } t]
  \label{eq:iv_v6}
\end{equation}
The instrument exploits aggregate FL supply variation at other
purchasing units in the same product market and year. We treat
$\hat{\beta}_{\text{IV}}$ as a measurement-error diagnostic rather
than a primary estimator: the leave-one-out instrument may correlate
with observable market characteristics that affect prices
independently of FL-classification noise, raising
exclusion-restriction concerns, and the point estimate is sensitive
to the FE structure. The IV indicates that the OLS is conservative,
but we do not include it in the headline range. Our primary
\emph{conditional-association} estimates are cross-fit (0.036), OLS
(0.064), and matching (0.055--0.077).

The conditional-association range above is the \emph{baseline}
estimate. The institutional identification of
Section~\ref{sec:empirical} above---modal comparison, within-convite
constraint binding, and the statutory value threshold---tightens
the reading by isolating the regulatory cover-bidding channel.
Section~\ref{sec:additional_id} reports the modal and constraint
estimates; together with the baseline conditional associations,
they bound the range over which the FL--price gap is concentrated
in the procurement environments where the regulatory incentive is
operative.

% ── 5.2 Detection Validation ─────────────────────────────────────

\subsection{Detection Performance Validation}
\label{sec:detection_validation}

The detection assessment uses CADE cartel adjudications as ground
truth. FL status is defined on 2009--2019 BEC participation, and
ground truth is co-bidding with the 47 BEC-active firm-defendants
across 12 cartel cases adjudicated 2015--2025. As discussed in
Section~\ref{sec:cade_sample}, the validation is partly prospective:
eight of the twelve cases were decided after our 2009--2019 sample
window, so the screen is asked to anticipate firms whose cartel
involvement subsequent enforcement confirmed.

\paragraph{ROC analysis.}
We vary the IQR multiplier from 0.5 to 5.0 in steps of 0.1, producing
46 different FL thresholds. For each threshold, we compute:
\begin{align}
  \text{TPR}(m) &= \frac{|\{j : \text{FL}_m(j)=1\} \cap
    \{j : \text{CADE}(j)=1\}|}{|\{j : \text{CADE}(j)=1\}|}
  \label{eq:tpr} \\
  \text{FPR}(m) &= \frac{|\{j : \text{FL}_m(j)=1\} \cap
    \{j : \text{CADE}(j)=0\}|}{|\{j : \text{CADE}(j)=0\}|}
  \label{eq:fpr}
\end{align}
where $\text{FL}_m(j)$ indicates FL classification under multiplier $m$
and $\text{CADE}(j)$ indicates co-participation with CADE-convicted
cartelists. AUC~$= 0.5$ corresponds to random classification and
AUC~$= 1.0$ to perfect discrimination.

\paragraph{Optimal threshold.}
We determine the optimal IQR multiplier using Youden's $J$
statistic:
\begin{equation}
  m^*_{\text{optimal}} = \operatorname*{arg\,max}_m \;
  \left[\text{TPR}(m) - \text{FPR}(m)\right]
  \label{eq:youden}
\end{equation}
which balances true-positive and false-positive rates, connecting the
IQR rule to a formal detection-theoretic criterion.

\paragraph{Benchmarking against Imhof-style screens.}
The FL screen is benchmarked against \citet{imhof2019detecting}-style
bid-level features (within-tender CV, kurtosis, skewness, normalized
spread) on the same CADE ground truth. A horse-race regression
including both FL presence and an Imhof-style CV flag in
Equation~\eqref{eq:ols_main} tests whether FL captures information
orthogonal to bid-level screens
(Section~\ref{sec:results_detection}).

% ── 5.3 Supporting diagnostics (online setup) ────────────────────

\subsection{Supporting Diagnostics: Setup}
\label{sec:supporting_setup}

Two supporting diagnostics complement the institutional
identification and the conditional-association baselines:
\textit{(i)~A two-stage structural mixture model} estimating
genuine-bid and cover-bid distributions on the bid-level data
(23.2 million non-FL losing bids and \valItemsCount FL losing bids);
the model selects between complementary and coordinated cover-bid
regimes via BIC and computes an implied markup as a calibration
exercise---not as an independently identified structural
parameter. \textit{(ii)~The \citet{bajari2003deciding}
bid-coordination tests} using corrected first-stage residuals
(excluding $n_{\text{bids}}$ as in \citealt{bajari2003deciding},
p.~978) to test exchangeability of FL and non-FL bid residuals
(KS test) and conditional independence of FL residuals within
tenders (mean pairwise product, bootstrap with \valPermPilot iterations,
plus a fake-groups placebo and a tender-FE variant). Full setup,
estimating equations, and identification discussion appear in
\ref{app:structural_id} (Online); structural results appear in
the same online appendix, and Bajari--Ye results are reported in
Section~\ref{sec:results_bajari}.

% ── 5.5 Network-Split Heterogeneity ──────────────────────────────

\subsection{Network-Split Heterogeneity Test}
\label{sec:empirical_network}

To test whether cover bidding concentrates in competitive markets
(a standard prediction we develop formally in
\ref{app:proofs} (Online)), we split FL firms into two
subgroups based on co-bidding network characteristics.

\paragraph{Winner HHI.}
For each FL firm $j$, we compute the Herfindahl--Hirschman Index of
winning firms across all tenders in which $j$ participates:
\begin{equation}
  \text{HHI}_j = \sum_{w=1}^{W_j} s_{wj}^2
  \label{eq:winner_hhi}
\end{equation}
where $s_{wj}$ is the share of firm $j$'s tenders won by winner $w$,
and $W_j$ is the number of distinct winners. A high $\text{HHI}_j$
indicates that $j$ operates in markets dominated by a single or few
winners (concentrated markets); a low $\text{HHI}_j$ indicates
diversified winner patterns (competitive markets).

\paragraph{Classification.}
We classify FL firms as \emph{concentrated-market} if
$\text{HHI}_j$ exceeds the sample median \emph{and} the firm shares
at least two tenders with at least one repeat partner (a necessary
condition for sustained coordination across tenders); and
\emph{competitive-market} as the complement. We then estimate the
baseline price regression (Equation~\ref{eq:ols_main}) separately
for each subgroup, interacting $\text{losers}_{igt}$ with the
market-structure indicator.

Before presenting the price regressions, we check the FL
classification against external enforcement data
(Section~\ref{sec:cade}). This provides an initial check that FL presence flags
economically meaningful environments---a necessary condition for the
price association to be informative.
